\section{总结}\label{sec:summary}

本文研究了以 $\gamma^t$ 定义的准 PDF，它在 $\mathcal{O}(a^0)$ 阶无算符混合。我们利用 $M_\pi=356$~MeV 的格点数据演示了匹配流程，并表明激发态污染得到了良好控制。我们计算了一圈匹配系数，并详细讨论了系统误差来源以及投影的选取。

我们发现来自 FT 截断方法和 $p_z^R$ 依赖的系统不确定度相当可观。但在 $p_z^R=0$ 时，来自 $\mu_R$、匹配反演以及投影选取的不确定度相对较小。与此同时，从准 PDF 到匹配 PDF 的显著变化表明需要更高圈阶的修正，如图~\ref{fig:matchingPDF} 所示。

控制激发态带来的系统不确定性极具挑战性，因为当源—汇分离 $t_\text{seq}$ 或核子动量 $P_z$ 变大时相对不确定度增长极快。较小分离下的二态拟合提供了在 $t_\text{seq}<1$ fm 小分离区获得精确结果的可能，而要在大分离处利用极高统计量作精确测量，仍需对此类拟合的系统不确定度作出估计。

除已研究的各类不确定度外，我们计划在未来考察其他系统误差，如格点离散化与有限体积效应~\cite{Lin:2019ocg}，以及影响小 $x$ 结果的更高扭度贡献。后者可通过更大的核子动量加以改善，并通过向无穷大核子动量外推进行估计。

我们关于光锥 PDF 的最终结果在大 $x$ 区与全局分析一致，这是一个令人鼓舞的信号，表明 LaMET 或许能使我们在未来精确地触及部分子物理。
